\chapter{Traffic-engineering methodology}
\label{chap:methodology}

\section{Method provenance}
Table~\ref{tab:method-provenance} distinguishes established engineering relationships
from PySig-specific implementations and approximations.

\begin{longtable}{p{0.29\textwidth}p{0.64\textwidth}}
\caption{Method provenance}\label{tab:method-provenance}\\
\toprule
\textbf{Quantity or process} & \textbf{Current PySig basis}\\
\midrule\endhead
Saturation flow & Direct measured/calibrated input or RR67-style unopposed estimate.\\
Turning proportion & Fixed override where supplied; otherwise derived from assigned lane movement flows.\\
Effective green & Displayed signal green minus start displacement plus end displacement, including continuous signal-group overlap treatment.\\
Nominal capacity & $c_n=s(g/C)$.\\
Nominal degree of saturation & $x_n=q/c_n$.\\
Simulation queue & Deterministic aggregate time-step conservation model.\\
Analytical queue & Deterministic event-driven fluid queue.\\
Queue-area delay & Integral of queue over time; trapezoidal exact integration for linear analytical segments.\\
Shared lane, common group & One physical lane queue/service, allocated across movement composition.\\
Different-group shared lane & Expected-value front-blocking approximation where configured/supported.\\
Priority movement & Explicit active capacity-factor approximation within physical lane supply.\\
Short/nested lanes & Lane-access tree using actual lane lengths, storage-limited per-cycle short-lane contribution and common upstream capacity.\\
Lane assignment & Logit target route choice with MSA flow update.\\
Signal optimisation & Sequential candidate evaluation around the supplied controller plan.\\
LOS & Named delay-based classification only when its delay basis is explicitly identified; queueing-delay screening must not be mislabelled as formal HCM control-delay LOS.\\
\bottomrule
\end{longtable}

\section{Fundamental queue relationships}
Traffic-flow theory at signals distinguishes arrivals, saturation discharge, effective
green, queue and delay.  PySig follows this decomposition.  The analytical implementation
assumes piecewise-constant rates between events.  Let $a$ be arrival rate and $\mu$ the
currently accessible service rate in PCU/s.  While $Q>0$,
\begin{equation}
\dot{Q}=a-\mu.
\end{equation}
For an interval $\tau$ that does not contain a queue-clearance event,
\begin{equation}
Q_1=Q_0+(a-\mu)\tau.
\end{equation}
The queue-area delay for a linear segment is
\begin{equation}
D=\frac{Q_0+Q_1}{2}\tau.
\end{equation}
This yields PCU-seconds.  Division by 3600 yields PCU-hours.

\section{Mean delay denominator}
Where PySig reports queueing delay per departed PCU,
\begin{equation}
\bar d_D=\frac{D}{N_D}.
\end{equation}
A model with residual queue may require additional care because dividing by departures
do not give the same interpretation as dividing by arrivals.  Outputs and LOS
classification must therefore identify the delay basis.

\section{Queue length in metres}
The simple reported spatial estimate is
\begin{equation}
L_Q=Q\,l_{PCU},
\end{equation}
where $l_{PCU}$ is configured effective PCU spacing.  This is not a kinematic shockwave
model.  ARR~340 provides empirical material on jam spacing and queue-clearance waves
that may support a future dynamic spatial-queue implementation.

\section{Australian empirical context}
Ak\c{c}elik, Besley and Roper's ARR~340 study used observations at 18 signalised
intersections in Sydney and Melbourne and examined queue-discharge headways and speeds,
jam spacing, saturation flow, start loss, end gain, departure response and
queue-clearance wave speed.  PySig uses this work as an important Australian validation
reference.  The present constant-saturation analytical service assumption should not be
confused with ARR~340's more detailed exponential discharge relationships.

\section{Calibration and validation}
A transparent formula is not a substitute for calibration.  Important models should be
checked against observed demand, lane use, saturation discharge, start/end effects,
queues and delays.  Field observations should be kept distinct from model calibration
parameters and from software defaults.
